%%% FIELD construction-clipped-ht
\(h_\alpha=\sqrt{S\log(2/\alpha)}+\frac{M}{6}\log(2/\alpha),\qquad I_{\mathrm{HT},\alpha}(D)=[-1,1]\cap[\widehat\theta(D)-h_\alpha,\widehat\theta(D)+h_\alpha]\).
%%% FIELD construction-trimmed-radius
\(\rho_\alpha(w,q)=\inf_{J\subseteq\{1,\ldots,L\}}[\sum_{l\notin J}w_l+\sqrt{\log(2/\alpha)\sum_{l\in J}w_l^2/q_l}+\{\log(2/\alpha)/6\}\max_{l\in J}w_l/q_l]\), with the maximum over \(J=\varnothing\) equal to zero, and \(\rho_0(w,q)=\inf_J[\sum_{l\notin J}w_l+(\sum_{l\in J}w_l^2/q_l)^{1/2}]\).
%%% FIELD statement-block-reduction
For every \(m\in\mathcal M_n(B)\), the vectors \((A_l,C_l,H_l)\) are independent across \(l\), each has categorical probabilities \((q_l,q_l,1-2q_l)\), and \(A_lC_l=0\). The observed-data construction in \(\mathrm{def{:}block\text{-}ht}\) satisfies \(\mathbb E_Z[X_l(D)]=w_l\{\mu_l(a)-\mu_l(b)\}\), \(\sum_l\operatorname{Var}_Z\{X_l(D)\}\le S/2\), and \(|X_l(D)-\mathbb E_ZX_l(D)|\le w_l/(2q_l)\le M/2\).
%%% FIELD statement-finite-upper
For every \(0<\alpha\le1/4\) and every \(m\in\mathcal M_n(B)\), the interval \(I_{\mathrm{HT},\alpha}(D)\), which is a measurable function solely of \(D\) and known design quantities, obeys \(P_{Z,m}\{\theta\in I_{\mathrm{HT},\alpha}(D)\}\ge1-\alpha\) and \(\operatorname{len}\{I_{\mathrm{HT},\alpha}(D)\}\le2\min\{1,\sqrt{S\log(2/\alpha)}+(M/6)\log(2/\alpha)\}\).
%%% FIELD statement-capped-trimming-duality
Let \(0<\alpha\le1/4\), \(x=\log(1/\alpha)\), and \(y=\log(2/\alpha)\). The optimizer in \(\mathrm{def{:}capped\text{-}modulus}\) is the stated water-filling vector. Moreover,
\[
 \omega_\alpha(w,q)\le\rho_\alpha(w,q)\le\frac52\omega_\alpha(w,q),
 \qquad
 \omega_1(w,q)\le\rho_0(w,q)\le2\omega_1(w,q),
\]
where \(\omega_1\) denotes the same capped program with quadratic budget one.
%%% FIELD statement-two-scale-rademacher-lower
For the family \(B^{(k)}\) containing \(k\) blocks of size two and \(k\) blocks of size four, every \(0<\alpha\le1/4\) and every integer \(k\ge64\) satisfy
\[
 \mathcal R_\alpha(w,q)\ge\frac{1}{192\sqrt{k}}.
\]
This lower bound applies to the full class of measurable honest interval procedures, including procedures externally randomized by \(R\).
%%% FIELD statement-unequal-two-scale-family
Let \(B^{(k)}\) contain \(k\) blocks of size two and \(k\) blocks of size four, so \(n=6k\). The two strata have \((w_l,q_l)=(1/(3k),1/4)\) and \((w_l,q_l)=(2/(3k),1/16)\), respectively, and
\[
 S=\frac{68}{9k},\qquad M=\frac{32}{3k},\qquad
 \frac{M\sqrt{\log(1/\alpha)}}{\sqrt S}
 =\frac{32\sqrt{\log(1/\alpha)}}{\sqrt{68k}}.
\]
Thus, for every fixed \(0<\alpha\le1/4\), the no-dominance condition holds for all sufficiently large \(k\). With \(C_1\) denoting the universal constant in \(\mathrm{thm{:}no\text{-}dominance\text{-}upper}\), for every sufficiently large \(k\),
\[
 \frac{1}{192\sqrt{k}}
 \le \mathcal R_\alpha(w,q)
 \le C_1\sqrt{\frac{68\log(1/\alpha)}{9k}},
\]
and consequently \(\mathcal R_\alpha(w,q)=\Theta_\alpha(k^{-1/2})\). The two water-filling breakpoints \(q_l/w_l\) are \(3k/4\) and \(3k/32\). No common block size can reproduce both reveal probabilities, or both breakpoints.
%%% FIELD statement-capped-modulus-upper
For every \(0<\alpha\le1/4\),
\[
 \mathcal R_\alpha(w,q)\le5\min\{1,\omega_\alpha(w,q)\}.
\]
The bound is attained by a deterministic observed-data interval and therefore remains valid when the competing class permits randomization by the independent external randomizer \(R\). In addition, if \(x_\alpha<\sum_{l=1}^Lq_l\) and no coordinate of the water-filling optimizer is capped, meaning \(\lambda_\alpha w_l/q_l<1\) for every \(l\), then
\[
 \omega_\alpha(w,q)=\sqrt{Sx_\alpha}=\sqrt{S\log(1/\alpha)}.
\]
%%% FIELD statement-water-filling-residual
Determine whether there is a universal constant \(c>0\) such that, for every \(0<\alpha\le1/4\) and every unequal-clique design in the model class,
\[
 \mathcal R_\alpha(w,q)\ge c\min\{1,\omega_\alpha(w,q)\}.
\]
The lower bound must hold when the infimum defining \(\mathcal R_\alpha\) ranges over every measurable interval procedure, including procedures randomized by the independent external randomizer \(R\). This universal linear lower bound is the sole residual question.
%%% FIELD water-filling-partial-result
The capped optimizer and both computable sides are closed. The strengthened trimming duality gives
\[
 \omega_\alpha(w,q)\le\rho_\alpha(w,q)\le\frac52\omega_\alpha(w,q),
 \qquad
 \omega_1(w,q)\le\rho_0(w,q)\le2\omega_1(w,q).
\]
Together with the trimmed interval, \(\mathrm{thm{:}capped\text{-}modulus\text{-}upper}\) proves
\[
 \mathcal R_\alpha(w,q)\le5\min\{1,\omega_\alpha(w,q)\}.
\]
When no cap binds, the active quadratic constraint gives \(\omega_\alpha(w,q)=\sqrt{S\log(1/\alpha)}\). On the converse side, \(\mathrm{lem{:}capped\text{-}erasure\text{-}consistency\text{-}lower}\) proves
\[
 \mathcal R_\alpha(w,q)
 \ge\frac{\omega_\alpha(w,q)^3}{128\log(1/\alpha)+8\omega_\alpha(w,q)^2}
 \ge\frac{\omega_\alpha(w,q)^3}{136\log(1/\alpha)}.
\]
Thus the only unresolved step is to replace this cubic converse by a universal linear converse at nonzero finite confidence.
%%% FIELD proof-block-reduction
Fix a block \(B_l\). By \(\mathrm{ass{:}balanced\text{-}bernoulli}\), each of its \(2^{r_l}\) assignment vectors has probability \(2^{-r_l}=q_l\). Exactly one vector is all treated and exactly one is all controlled. Thus
\[
 P_Z(A_l=1)=q_l,\qquad P_Z(C_l=1)=q_l,
 \qquad P_Z(H_l=1)=1-2q_l,
\tag{1}
\]
and \(A_lC_l=0\) on every assignment. By \(\mathrm{ass{:}partition}\), distinct triples \((A_l,C_l,H_l)\) depend on disjoint coordinates of \(Z\); their independence follows from \(\mathrm{ass{:}balanced\text{-}bernoulli}\).

Put \(u_l=\mu_l(a)\) and \(v_l=\mu_l(b)\). The fixed-schedule and bounded-outcome assumptions, together with \(\mathrm{def{:}block\text{-}means}\), give
\[
 -\frac12\le u_l,v_l\le\frac12.
\tag{2}
\]
On the all-treated and all-control events, \(\mathrm{lem{:}observed\text{-}block\text{-}means}\) identifies the observed block mean with \(u_l\) and \(v_l\), respectively. Mutual exclusion then rewrites the observed statistic, on every assignment, as
\[
 X_l=\frac{w_l}{q_l}(A_lu_l-C_lv_l).
\tag{3}
\]
Using (1), fixedness of \((u_l,v_l)\), and \(A_lC_l=0\),
\[
 \mathbb E_ZX_l=w_l(u_l-v_l),
\tag{4}
\]
and
\[
\begin{aligned}
 \operatorname{Var}_Z(X_l)
 &\le\mathbb E_ZX_l^2\\
 &=\frac{w_l^2}{q_l}(u_l^2+v_l^2)\\
 &\le\frac{w_l^2}{2q_l}.
\end{aligned}
\tag{5}
\]
The last step uses (2). Summing (5) and using the definition of \(S\) proves
\[
 \sum_{l=1}^L\operatorname{Var}_Z(X_l)
 \le\frac12\sum_{l=1}^L\frac{w_l^2}{q_l}=\frac S2.
\tag{6}
\]

It remains to compute the exact centered increment. Write \(W_l=X_l-\mathbb E_ZX_l\). Equations (3)--(4) give the following three exhaustive cases:
\[
 \frac{W_l}{w_l}=
 \begin{cases}
  (q_l^{-1}-1)u_l+v_l,&A_l=1,\\
  -u_l-(q_l^{-1}-1)v_l,&C_l=1,\\
  -u_l+v_l,&H_l=1.
 \end{cases}
\tag{7}
\]
Since \(r_l\ge1\), one has \(0<q_l\le1/2\). By (2), the first two rows of (7) have absolute value at most
\[
 \frac12(q_l^{-1}-1)+\frac12=\frac1{2q_l}.
\tag{8}
\]
The third row has absolute value at most one, which is at most \(1/(2q_l)\) because \(q_l\le1/2\). Hence on every assignment
\[
 |X_l-\mathbb E_ZX_l|=|W_l|
 \le\frac{w_l}{2q_l}\le\frac M2.
\tag{9}
\]
This is an enumeration of All Treat, All Control, and every other assignment, rather than a triangle bound on \(|X_l|+|\mathbb E_ZX_l|\).

For the smallest unequal stress test, take a singleton block and a two-unit block. Their \(2\cdot4=8\) joint assignments give categorical counts \((1,1,0)\) and \((1,1,2)\), respectively. Substitution in (7), for arbitrary \(u_l,v_l\in[-1/2,1/2]\), verifies (8)--(9) on each of the eight assignments. Thus the asserted support, mean, variance bound, and increment bound also hold exactly at this finite boundary case.
%%% FIELD proof-finite-upper
Fix \(m\in\mathcal M_n(B)\), and define
\[
 W_l=X_l-\mathbb E_ZX_l,\qquad z=\log(2/\alpha)>0.
\tag{10}
\]
By \(\mathrm{lem{:}block\text{-}reduction}\), the variables \(W_l\) are independent and centered and obey
\[
 v:=\sum_{l=1}^L\mathbb E_ZW_l^2\le\frac S2,
 \qquad |W_l|\le\frac M2.
\tag{11}
\]
Apply \(\mathrm{lem{:}finite\text{-}product\text{-}exponential}\) with \(c=M/2\) and the value of \(z\) in (10). Because its threshold is nondecreasing in \(v\), (11) gives
\[
\begin{aligned}
 P_Z\!\left\{\left|\sum_{l=1}^LW_l\right|>
 \sqrt{S\log(2/\alpha)}+\frac M6\log(2/\alpha)\right\}
 &\le2\exp\{-\log(2/\alpha)\}\\
 &=\alpha.
\end{aligned}
\tag{12}
\]
The expectation identity in \(\mathrm{lem{:}block\text{-}reduction}\), followed by the definition of the causal estimand, yields
\[
 \sum_{l=1}^L\mathbb E_ZX_l
 =\sum_{l=1}^Lw_l\{\mu_l(a)-\mu_l(b)\}=\theta.
\tag{13}
\]
Since \(\widehat\theta=\sum_lX_l\) by \(\mathrm{def{:}block\text{-}ht}\), equations (12)--(13) and the corrected construction in \(\mathrm{def{:}clipped\text{-}ht}\) imply
\[
 P_{Z,m}\{|\widehat\theta-\theta|\le h_\alpha\}\ge1-\alpha.
\tag{14}
\]
Every \(X_l\), and hence \(\widehat\theta\), is measurable from \(D\) by \(\mathrm{def{:}block\text{-}ht}\) and \(\mathrm{lem{:}observed\text{-}block\text{-}means}\). The quantities \(S\), \(M\), and \(\alpha\) are known design inputs. Moreover, bounded outcomes give \(\mu_l(a)-\mu_l(b)\in[-1,1]\), and the partition together with \(w_l=r_l/n\) gives \(\sum_lw_l=1\). Thus \(\theta\in[-1,1]\). On the event in (14), intersecting the raw interval with \([-1,1]\) cannot remove \(\theta\), proving
\[
 P_{Z,m}\{\theta\in I_{\mathrm{HT},\alpha}(D)\}\ge1-\alpha.
\tag{15}
\]
Finally, the length of an intersection is at most the length of each intersected interval, whence, pointwise in \(D\),
\[
\begin{aligned}
 \operatorname{len}\{I_{\mathrm{HT},\alpha}(D)\}
 &\le\min\{2,2h_\alpha\}\\
 &=2\min\!\left\{1,\sqrt{S\log(2/\alpha)}
       +\frac M6\log(2/\alpha)\right\}.
\end{aligned}
\tag{16}
\]
This deterministic finite-sample certificate holds on every assignment, including all eight assignments in the unequal three-unit stress test of \(\mathrm{lem{:}block\text{-}reduction}\).
%%% FIELD proof-trimmed-upper
Fix a deterministic subset \(J\subseteq\{1,\ldots,L\}\) and define
\[
 \widehat\theta_J=\sum_{l\in J}X_l,
 \qquad
 \theta_J=\sum_{l\in J}w_l\{\mu_l(a)-\mu_l(b)\},
\]
\[
 S_J=\sum_{l\in J}\frac{w_l^2}{q_l},\qquad
 M_J=\max_{l\in J}\frac{w_l}{q_l},\qquad
 b_J=\sum_{l\notin J}w_l,
\tag{17}
\]
with \(M_\varnothing=0\). For \(l\in J\), \(\mathrm{lem{:}block\text{-}reduction}\) gives independent centered variables \(W_l=X_l-\mathbb E_ZX_l\) satisfying
\[
 \sum_{l\in J}\mathbb E_ZW_l^2\le\frac{S_J}{2},
 \qquad
 |W_l|\le\frac{w_l}{2q_l}\le\frac{M_J}{2}.
\tag{18}
\]
If \(J\ne\varnothing\), apply \(\mathrm{lem{:}finite\text{-}product\text{-}exponential}\) with \(v\le S_J/2\), \(c=M_J/2\), and \(y=\log(2/\alpha)>0\). If \(J=\varnothing\), the same conclusion holds deterministically because both sums are zero. In either case,
\[
 P_Z\!\left\{|\widehat\theta_J-\theta_J|>
 \sqrt{S_Jy}+\frac{M_Jy}{6}\right\}\le\alpha.
\tag{19}
\]
Bounded outcomes imply \(|\mu_l(a)-\mu_l(b)|\le1\), so
\[
 |\theta-\theta_J|
 =\left|\sum_{l\notin J}w_l\{\mu_l(a)-\mu_l(b)\}\right|
 \le b_J.
\tag{20}
\]
Consequently the observed-data interval
\[
 I_J(D)=[-1,1]\cap
 \left[\widehat\theta_J-b_J-\sqrt{S_Jy}-\frac{M_Jy}{6},
       \widehat\theta_J+b_J+\sqrt{S_Jy}+\frac{M_Jy}{6}\right]
\tag{21}
\]
has coverage at least \(1-\alpha\) by (19)--(20). Each \(X_l\) is measurable from \(D\), and \(J\), \(S_J\), \(M_J\), and \(b_J\) are fixed known-design quantities, so \(I_J\) is measurable from \(D\). Its pointwise length is at most
\[
 2\min\!\left\{1,b_J+\sqrt{S_Jy}+\frac{M_Jy}{6}\right\}.
\tag{22}
\]
There are finitely many subsets. Choose a minimizer of the known expression in (22), using a fixed deterministic tie rule. The selected subset is nonrandom, so (19) retains its coverage. By the corrected construction in \(\mathrm{def{:}trimmed\text{-}radius}\), the expression after the leading one in (22) is exactly \(\rho_\alpha(w,q)\). Thus the selected interval lies in \(\mathfrak C_\alpha\) and has length at most \(2\min\{1,\rho_\alpha(w,q)\}\), as claimed.
%%% FIELD proof-capped-trimming-duality
Because every block is nonempty, \(w_l>0\); and because \(q_l=2^{-r_l}\), also \(q_l>0\). For \(z>0\), define
\[
 \mathcal D_z=\left\{\delta\in[0,1]^L:
                 \sum_{l=1}^Lq_l\delta_l^2\le z\right\}.
\tag{23}
\]
This is a nonempty compact subset of \(\mathbb R^L\), and the linear objective is continuous, so a maximizer exists. If \(z\ge\sum_lq_l\), then \((1,\ldots,1)\in\mathcal D_z\); it is optimal because \(w_l>0\) and every feasible coordinate is at most one.

Suppose \(0<z<\sum_lq_l\). The function
\[
 f(\lambda)=\sum_{l=1}^Lq_l\min\{1,\lambda w_l/q_l\}^2,
 \qquad \lambda\ge0,
\tag{24}
\]
is continuous, equals zero at zero, and is strictly increasing until all coordinates reach their caps, after which it equals \(\sum_lq_l\). Hence there is a unique \(\lambda>0\), before the terminal plateau, such that \(f(\lambda)=z\). Put
\[
 \delta_l^*=\min\{1,\lambda w_l/q_l\}.
\tag{25}
\]
This vector belongs to \(\mathcal D_z\) and makes its quadratic constraint active.

For optimality, let \(C=\{l:\delta_l^*=1\}\), let \(J=C^c\), and take any \(\delta\in\mathcal D_z\). For \(l\in J\), the identity \(w_l=q_l\delta_l^*/\lambda\) and \((\delta_l-\delta_l^*)^2\ge0\) give
\[
 w_l(\delta_l-\delta_l^*)
 \le\frac{q_l}{2\lambda}\{\delta_l^2-(\delta_l^*)^2\}.
\tag{26}
\]
For \(l\in C\), the cap relation gives \(w_l\ge q_l/\lambda\). Since \(\delta_l-1\le0\) and \((\delta_l-1)^2\ge0\),
\[
 w_l(\delta_l-1)
 \le\frac{q_l}{\lambda}(\delta_l-1)
 \le\frac{q_l}{2\lambda}(\delta_l^2-1).
\tag{27}
\]
Summing (26)--(27), using feasibility of \(\delta\), and using the active constraint for \(\delta^*\), yields
\[
 \sum_{l=1}^Lw_l(\delta_l-\delta_l^*)
 \le\frac1{2\lambda}\sum_{l=1}^Lq_l\{\delta_l^2-(\delta_l^*)^2\}
 \le0.
\tag{28}
\]
Thus (25) is the global optimizer. Equations (26)--(28) are also a first-order certificate for the capped convex feasible program.

We next compare the moduli. Fix \(\delta\) feasible at budget \(x\), and fix \(J_0\subseteq\{1,\ldots,L\}\). Since \(0\le\delta_l\le1\), Cauchy--Schwarz gives
\[
\begin{aligned}
 \sum_{l=1}^Lw_l\delta_l
 &\le\sum_{l\notin J_0}w_l+
      \sum_{l\in J_0}w_l\delta_l\\
 &\le\sum_{l\notin J_0}w_l+
 \left(\sum_{l\in J_0}\frac{w_l^2}{q_l}\right)^{1/2}
 \left(\sum_{l\in J_0}q_l\delta_l^2\right)^{1/2}\\
 &\le\sum_{l\notin J_0}w_l+
 \sqrt{x\sum_{l\in J_0}w_l^2/q_l}.
\end{aligned}
\tag{29}
\]
The corrected construction of \(\rho_\alpha\) has square-root coefficient \(\sqrt y\), where \(y=x+\log2\ge x\), and a nonnegative linear term. Taking the supremum over \(\delta\) and then the infimum over \(J_0\) in (29) proves
\[
 \omega_\alpha(w,q)\le\rho_\alpha(w,q).
\tag{30}
\]
If \(x\ge\sum_lq_l\), the optimizer is all-capped and \(\omega_\alpha=\sum_lw_l\). Choosing \(J_0=\varnothing\) gives \(\rho_\alpha\le\sum_lw_l=\omega_\alpha\). Now suppose \(x<\sum_lq_l\). Apply (25) at budget \(x\), retain its capped set \(C\) and uncapped set \(J\), and define
\[
 A=\sum_{l\in J}\frac{w_l^2}{q_l},\qquad
 Q_C=\sum_{l\in C}q_l,\qquad
 B=\sum_{l\in C}w_l.
\tag{31}
\]
The active constraint and objective value are
\[
 x=Q_C+\lambda^2A,
 \qquad
 \omega_\alpha=B+\lambda A.
\tag{32}
\]
For \(l\in C\), \(\lambda w_l/q_l\ge1\), so \(B\ge Q_C/\lambda\). For \(l\in J\), \(w_l/q_l<1/\lambda\). Consequently (32) implies
\[
 \omega_\alpha\ge\frac{x}{\lambda},
 \qquad A\le\frac{x}{\lambda^2},
 \qquad \max_{l\in J}\frac{w_l}{q_l}<\frac1\lambda.
\tag{33}
\]
Because \(0<\alpha\le1/4\), \(x\ge\log4=2\log2\), and hence \(y=x+\log2\le3x/2\). Choosing \(J_0=J\) in the corrected definition, whose linear coefficient is \(y/6\), and applying (32)--(33), gives
\[
\begin{aligned}
 \rho_\alpha
 &\le B+\sqrt{yA}+\frac y6\max_{l\in J}\frac{w_l}{q_l}\\
 &\le B+\sqrt{\frac32}\,\frac{x}{\lambda}
       +\frac14\frac{x}{\lambda}\\
 &\le\left(1+\sqrt{\frac32}+\frac14\right)\omega_\alpha
 <\frac52\omega_\alpha.
\end{aligned}
\tag{34}
\]
This proves the confidence-level comparison.

Finally, repeat (29) with quadratic budget one and without the linear term to obtain \(\omega_1\le\rho_0\). If \(\sum_lq_l\le1\), the budget-one optimizer is all-capped, so choosing the empty retained set gives \(\rho_0\le\sum_lw_l=\omega_1\). If \(\sum_lq_l>1\), apply (31)--(33) with \(x=1\). Choosing the uncapped set in the definition of \(\rho_0\) gives
\[
 \rho_0\le B+\sqrt A
 \le B+\lambda^{-1}
 \le2\omega_1.
\tag{35}
\]
Thus \(\omega_1\le\rho_0\le2\omega_1\), completing the proof.
%%% FIELD proof-capped-modulus-upper
Fix \(0<\alpha\le1/4\). By \(\mathrm{thm{:}trimmed\text{-}upper}\), there is a deterministic interval in \(\mathfrak C_\alpha\) whose pointwise length, under every fixed schedule in \(\mathcal M_n(B)\), is at most
\[
 2\min\{1,\rho_\alpha(w,q)\}.
\tag{36}
\]
Because the infimum in \(\mathrm{def{:}length\text{-}risk}\) ranges over a class containing this rule, taking the expectation, the supremum over schedules, and then the infimum over procedures gives
\[
 \mathcal R_\alpha(w,q)\le2\min\{1,\rho_\alpha(w,q)\}.
\tag{37}
\]
By \(\mathrm{lem{:}capped\text{-}trimming\text{-}duality}\), \(\rho_\alpha(w,q)\le(5/2)\omega_\alpha(w,q)\). For \(u\ge0\) and \(K\ge1\), \(\min\{1,Ku\}\le K\min\{1,u\}\), by considering \(u\le1\) and \(u>1\). Therefore
\[
 \mathcal R_\alpha(w,q)\le5\min\{1,\omega_\alpha(w,q)\}.
\tag{38}
\]
The witness is deterministic, so allowing competitors to use the independent \(R\) does not affect this upper bound.

For the no-cap identity, put \(x=x_\alpha\). Under \(x<\sum_lq_l\), the optimizer certificate gives
\[
 \delta_l^*=\min\{1,\lambda_\alpha w_l/q_l\},
 \qquad
 \sum_lq_l(\delta_l^*)^2=x.
\tag{39}
\]
If no cap binds, \(\delta_l^*=\lambda_\alpha w_l/q_l\) for every \(l\). Positivity of \(q_l\), \(w_l\), and \(S=\sum_lw_l^2/q_l\) permits division, and (39) becomes
\[
 x=\lambda_\alpha^2S,
 \qquad
 \lambda_\alpha=\sqrt{x/S}.
\tag{40}
\]
Substitution into the objective gives
\[
 \omega_\alpha(w,q)
 =\lambda_\alpha\sum_l\frac{w_l^2}{q_l}
 =\lambda_\alpha S
 =\sqrt{Sx}
 =\sqrt{S\log(1/\alpha)}.
\tag{41}
\]
%%% FIELD proof-no-dominance-upper
Put \(x=\log(1/\alpha)\) and \(y=\log(2/\alpha)\). Since \(0<\alpha\le1/4\), \(x\ge\log4\), and hence
\[
 y=x+\log2\le\frac32x.
\tag{42}
\]
The corrected \(\mathrm{thm{:}finite\text{-}upper}\) supplies a deterministic observed-data interval that is honest uniformly over \(\mathcal M_n(B)\) and has pointwise length at most
\[
 2\min\left\{1,\sqrt{Sy}+\frac{My}{6}\right\}.
\tag{43}
\]
It is an admissible deterministic member of \(\mathfrak C_\alpha\). Although \(\mathrm{def{:}length\text{-}risk}\) also permits procedures using \(R\), its infimum is bounded above by this rule. Thus (43) bounds \(\mathcal R_\alpha(w,q)\).

Take \(c_0=1\). Under \(M\sqrt{x}\le\sqrt S\), (42) gives
\[
\begin{aligned}
 \sqrt{Sy}+\frac{My}{6}
 &\le\sqrt{\frac32}\sqrt{Sx}+\frac14Mx\\
 &\le\left(\sqrt{\frac32}+\frac14\right)\sqrt{Sx}.
\end{aligned}
\tag{44}
\]
For \(K\ge1\) and \(t\ge0\), \(\min\{1,Kt\}\le K\min\{1,t\}\). Hence (43)--(44) prove the theorem with the universal constants
\[
 c_0=1,
 \qquad
 C_1=2\left(\sqrt{\frac32}+\frac14\right).
\tag{45}
\]
%%% FIELD proof-two-scale-rademacher-lower
Fix \(0<\alpha\le1/4\), an integer \(k\ge64\), and \(\Gamma\in\mathfrak C_\alpha\). Index the \(k\) size-two blocks by \(j=1,\ldots,k\). In the hard schedules from \(\mathrm{def{:}hard\text{-}subexperiment\text{-}embedding}\), put
\[
 V_j=\frac{\varepsilon_j}{2},\qquad
 P(\varepsilon_j=1)=P(\varepsilon_j=-1)=\frac12,
 \quad j=1,\ldots,k,
\tag{46}
\]
independently, and set \(V_l=1/2\) deterministically on every size-four block. This product law is used only to average the pointwise coverage and risk identities in \(\mathrm{lem{:}hard\text{-}subexperiment\text{-}embedding}\); every schedule in its support is a fixed member of \(\mathcal M_n(B^{(k)})\). For a size-two block,
\[
 q_j=\frac14,\qquad w_j=\frac1{3k}.
\tag{47}
\]
Let \(K_j\) be its all-treated reveal indicator and define \(N=\sum_{j=1}^k(1-K_j)\). The embedding lemma gives independent \(K_j\sim\operatorname{Bernoulli}(1/4)\), independent of the signs, and hence
\[
 N\sim\operatorname{Binomial}(k,3/4),\qquad
 \mathbb EN=\frac{3k}{4},\qquad
 \operatorname{Var}(N)=\frac{3k}{16}.
\tag{48}
\]
Write \(\Lambda=\Lambda_\Gamma(O,R)\) for the corresponding interval in the erasure experiment and \(L_\Lambda=\operatorname{len}(\Lambda)\). Conditional on \((O,R)\), the revealed size-two signs and the size-four contribution are fixed, whereas the \(N\) unrevealed size-two signs remain independent fair Rademacher variables. The independence of \(R\) is part of \(\mathrm{def{:}weighted\text{-}erasure\text{-}law}\). Since the embedding gives \(\theta=\sum_lw_lV_l\), the conditional target law is
\[
 c(O)+\frac1{6k}\sum_{s=1}^{N}\xi_s,
\tag{49}
\]
where the \(\xi_s\) are conditionally independent fair Rademacher variables. Adjacent support points in (49) are separated by \(1/(3k)\).

For \(S_N=\sum_{s=1}^N\xi_s\), its largest atom is \(p_N=2^{-N}\binom{N}{\lfloor N/2\rfloor}\). If \(N=2m\), factorial cancellation gives
\[
 p_{2m}=4^{-m}\binom{2m}{m}
       =\prod_{s=1}^m\left(1-\frac1{2s}\right).
\]
Since \(\log(1-u)\le-u\) for \(0<u<1\) and \(\sum_{s=1}^m s^{-1}\ge\int_1^{m+1}t^{-1}\,dt=\log(m+1)\),
\[
 p_{2m}\le(m+1)^{-1/2}\le\sqrt{\frac2{2m}}.
\tag{50}
\]
If \(N=2m+1\), then
\[
 p_{2m+1}=\frac{2m+1}{2m+2}p_{2m}
 \le(m+1)^{-1/2}
 \le\sqrt{\frac2{2m+1}}.
\tag{51}
\]
Thus for every \(N\ge1\), \(p_N\le\sqrt{2/N}\).

An interval of length \(u\) contains at most \(3ku+1\) points of a lattice with spacing \(1/(3k)\). Therefore, on \(N\ge k/2\), conditionally on \((O,R)\),
\[
 P\{\theta\in\Lambda\mid O,R\}
 \le(3kL_\Lambda+1)\sqrt{\frac2N}.
\tag{52}
\]
On the event
\[
 \mathcal A=\left\{N\ge\frac k2,\quad
 L_\Lambda<\frac1{48\sqrt{k}}\right\},
\]
equation (52) and \(k\ge64\) give
\[
 P\{\theta\in\Lambda\mid O,R\}
 \le\left(\frac{\sqrt{k}}{16}+1\right)\frac2{\sqrt{k}}
 =\frac18+\frac2{\sqrt{k}}
 \le\frac38.
\tag{53}
\]
By Chebyshev's inequality and (48),
\[
 P\left(N<\frac k2\right)
 \le\frac{3k/16}{(k/4)^2}
 =\frac3k\le\frac3{64}.
\tag{54}
\]
Pointwise honesty of \(\Gamma\), averaged over (46), and the coverage identity in the embedding lemma imply
\[
 P\{\theta\in\Lambda\}\ge1-\alpha\ge\frac34.
\tag{55}
\]
On \(\mathcal A\), conditional coverage is at most \(3/8\), and outside it is at most one. Since \(\mathcal A^c\) is contained in the union of \(\{N<k/2\}\) and \(\{L_\Lambda\ge1/(48\sqrt{k})\}\), equations (53)--(55) give
\[
 P\left(L_\Lambda\ge\frac1{48\sqrt{k}}\right)
 \ge\frac38-\frac3{64}=\frac{21}{64}\ge\frac14.
\tag{56}
\]
Therefore
\[
 \mathbb E L_\Lambda
 \ge\frac1{48\sqrt{k}}
 P\left(L_\Lambda\ge\frac1{48\sqrt{k}}\right)
 \ge\frac1{192\sqrt{k}}.
\tag{57}
\]
The expected-length identity and supremum inequality in the embedding lemma show that the worst fixed-schedule risk of the original \(\Gamma\) is at least the prior average in (57). Since \(\Gamma\) was arbitrary, taking the infimum in \(\mathrm{def{:}length\text{-}risk}\) proves the result. Conditioning throughout on the declared independent \(R\) covers every externally randomized measurable interval.
%%% FIELD proof-unequal-two-scale-family
The population size is \(2k+4k=6k\). Hence a size-two block has
\[
 w_l=\frac2{6k}=\frac1{3k},\qquad q_l=2^{-2}=\frac14,
\]
whereas a size-four block has
\[
 w_l=\frac4{6k}=\frac2{3k},\qquad q_l=2^{-4}=\frac1{16}.
\]
The reveal probabilities follow from \(\mathrm{ass{:}balanced\text{-}bernoulli}\). Substitution into \(S\) and \(M\) gives
\[
\begin{aligned}
 S
 &=k\frac{(1/(3k))^2}{1/4}
   +k\frac{(2/(3k))^2}{1/16}
   =\frac4{9k}+\frac{64}{9k}
   =\frac{68}{9k},\\
 M
 &=\max\left\{\frac{1/(3k)}{1/4},
                \frac{2/(3k)}{1/16}\right\}
   =\frac{32}{3k}.
\end{aligned}
\tag{58}
\]
Consequently
\[
 \frac{M\sqrt{\log(1/\alpha)}}{\sqrt S}
 =\frac{32\sqrt{\log(1/\alpha)}}{\sqrt{68k}}
 \longrightarrow0
\tag{59}
\]
for fixed \(\alpha\). Thus \(\mathrm{thm{:}no\text{-}dominance\text{-}upper}\) applies for all sufficiently large \(k\), and (58) yields
\[
 \mathcal R_\alpha(w,q)
 \le C_1\sqrt{\frac{68\log(1/\alpha)}{9k}}.
\tag{60}
\]
For \(k\ge64\), \(\mathrm{lem{:}two\text{-}scale\text{-}rademacher\text{-}lower}\) gives
\[
 \mathcal R_\alpha(w,q)\ge\frac1{192\sqrt{k}}.
\tag{61}
\]
Equations (60)--(61) prove \(\mathcal R_\alpha(w,q)=\Theta_\alpha(k^{-1/2})\). The lower bound is specialized to the repeated size-two stratum and makes no assertion about either universal weighted-erasure open question.

Finally, the cap breakpoints are \(q_l/w_l\), namely
\[
 \frac{1/4}{1/(3k)}=\frac{3k}{4},
 \qquad
 \frac{1/16}{2/(3k)}=\frac{3k}{32}.
\tag{62}
\]
Since \(r\mapsto2^{-r}\) is injective on the positive integers, one common block size cannot reproduce both reveal probabilities. The two distinct values in (62) likewise rule out a common breakpoint.
%%% FIELD proof-consistency-boundary
Fix \(0<\alpha\le1/4\), put \(x=\log(1/\alpha)>1\), and write
\[
 \omega_z(w,q)=\max\left\{\sum_{l=1}^Lw_l\delta_l:
 0\le\delta_l\le1,\quad
 \sum_{l=1}^Lq_l\delta_l^2\le z\right\}.
\tag{63}
\]
Thus \(\omega_\alpha=\omega_x\), whereas \(\omega_1\) is the budget-one modulus. Since \(x>1\), nesting gives \(\omega_1\le\omega_x\). Conversely, if \(\delta\) is feasible at budget \(x\), then \(\delta/\sqrt{x}\in[0,1]^L\), its quadratic cost is at most one, and its objective is the original objective divided by \(\sqrt{x}\). Taking suprema proves
\[
 \omega_1\le\omega_\alpha\le\sqrt{x}\,\omega_1.
\tag{64}
\]
By \(\mathrm{lem{:}capped\text{-}trimming\text{-}duality}\),
\[
 \omega_1\le\rho_0\le2\omega_1.
\tag{65}
\]
Because \(x\) is fixed along the triangular array, (64)--(65) imply
\[
 \omega_\alpha\to0\quad\Longleftrightarrow\quad\rho_0\to0.
\tag{66}
\]
The unrestricted upper theorem gives
\[
 \mathcal R_\alpha(w,q)\le5\min\{1,\omega_\alpha(w,q)\},
\tag{67}
\]
so \(\omega_\alpha\to0\) implies \(\mathcal R_\alpha\to0\). Conversely, \(\mathrm{lem{:}capped\text{-}erasure\text{-}consistency\text{-}lower}\) gives row by row
\[
 \mathcal R_\alpha(w,q)\ge\frac{\omega_\alpha(w,q)^3}{136x}.
\tag{68}
\]
Since \(136x\) is fixed, finite, and positive, \(\mathcal R_\alpha\to0\) forces \(\omega_\alpha\to0\). Equations (66)--(68) prove the three-way equivalence.

For sufficiency of \(S^{(n)}\to0\), positivity of every \(q_l\) and Cauchy--Schwarz give, for each vector feasible at budget \(x\),
\[
\begin{aligned}
 \sum_{l=1}^Lw_l\delta_l
 &=\sum_{l=1}^L\frac{w_l}{\sqrt{q_l}}\sqrt{q_l}\delta_l\\
 &\le\left(\sum_{l=1}^L\frac{w_l^2}{q_l}\right)^{1/2}
       \left(\sum_{l=1}^Lq_l\delta_l^2\right)^{1/2}
 \le\sqrt{xS}.
\end{aligned}
\tag{69}
\]
Maximizing proves \(\omega_\alpha\le\sqrt{xS}\); (67) then shows that \(S^{(n)}\to0\) is sufficient.

For nonnecessity, for each integer \(k\ge1\) set \(n_k=2k\,2^k\). Partition the \(n_k\) units into one block of size \(2k\) and \(n_k-2k\) singleton blocks. These blocks are nonempty, disjoint, and exhaustive. Under balanced Bernoulli assignment, the large block has
\[
 w_{\mathrm{big}}=\frac{2k}{n_k}=2^{-k},
 \qquad q_{\mathrm{big}}=2^{-2k},
 \qquad \frac{w_{\mathrm{big}}^2}{q_{\mathrm{big}}}=1.
\tag{70}
\]
Every singleton has \(w_l=1/n_k\) and \(q_l=1/2\). Therefore
\[
 S^{(n_k)}=1+\frac{2(n_k-2k)}{n_k^2}\longrightarrow1.
\tag{71}
\]
Let \(J_k\) contain all singleton blocks and put \(y=\log(2/\alpha)\). The omitted weight is \(2^{-k}\), the retained quadratic sum is \(2(n_k-2k)/n_k^2\), and the largest retained ratio is \(2/n_k\). The corrected trimmed radius gives
\[
 \rho_\alpha(w^{(n_k)},q^{(n_k)})
 \le2^{-k}+\sqrt{\frac{2y(n_k-2k)}{n_k^2}}
       +\frac{y}{3n_k}
 \longrightarrow0.
\tag{72}
\]
The lower comparison \(\omega_\alpha\le\rho_\alpha\), followed by (67), gives \(\mathcal R_\alpha(w^{(n_k)},q^{(n_k)})\to0\), whereas (71) shows \(S^{(n_k)}\not\to0\). This explicit unequal-clique triangular array proves nonnecessity.
